% Electronic supplementary material. Compile separately.
\documentclass[11pt,a4paper]{article}
\usepackage[T1]{fontenc}
\usepackage[utf8]{inputenc}
\usepackage{lmodern}
\usepackage[margin=2.5cm]{geometry}
\usepackage{graphicx}
\usepackage[labelfont=bf,font=small]{caption}
\usepackage{xcolor}
\renewcommand{\thefigure}{S\arabic{figure}}
\newcommand{\todo}[1]{\textcolor{red}{[#1]}}

\begin{document}
\begin{center}
{\large\bfseries Electronic supplementary material for ``\todo{Title}''}\\[0.4em]
\todo{Authors}
\end{center}

\begin{figure}[h]
\centering
\includegraphics[width=\textwidth]{figures/figS1_intrinsic_dynamics.pdf}
\caption{Dynamics of an isolated patch. (\emph{a}) Psyllid density without releases (grey) and with releases of $\delta = 35$ every $T = 15$ (blue), from the initial state $(50,10,15,30)$. (\emph{b}) Dispersion of the state sampled after each of the last 30 of 120 releases, on a $\log_{10}$ scale. Zero dispersion would indicate a response of period $T$. No schedule on the grid reached a dispersion below 3.2.}
\end{figure}

\begin{figure}[h]
\centering
\includegraphics[width=0.6\textwidth]{figures/figS2_sensitivity.pdf}
\caption{One-at-a-time sensitivity of the network pest AUC to multiplying each local parameter by 0.9 or 1.1. Complete graph, uniform releases, $\delta = 35$, $T = 15$, $D_d = D_t = 0.03$. Parameters are ordered by the larger absolute change. For negative parameters the factor 1.1 increases their magnitude.}
\end{figure}
\end{document}
